%%%%%%%%%%%%%%%%%%%%%%%%%%%%%%%%%%%%%%%%%%%%%%%%%%%%%%%%%%%%%%%%%%%%%%%%%%%
%
% Generic template for TFC/TFM/TFG/Tesis
%
% $Id: desarrollo.tex,v 1.4 2015/06/05 00:05:18 macias Exp $
%
% By:
%  + Javier Macías-Guarasa.
%    Departamento de Electrónica
%    Universidad de Alcalá
%  + Roberto Barra-Chicote.
%    Departamento de Ingeniería Electrónica
%    Universidad Politécnica de Madrid
% 
% Based on original sources by Roberto Barra, Manuel Ocaña, Jesús Nuevo, Pedro Revenga, Fernando Herránz and Noelia Hernández. Thanks a lot to all of them, and to the many anonymous contributors found (thanks to google) that provided help in setting all this up.
%
% See also the additionalContributors.txt file to check the name of additional contributors to this work.
%
% If you think you can add pieces of relevant/useful examples, improvements, please contact us at (macias@depeca.uah.es)
%
% You can freely use this template and please contribute with comments or suggestions!!!
%
%%%%%%%%%%%%%%%%%%%%%%%%%%%%%%%%%%%%%%%%%%%%%%%%%%%%%%%%%%%%%%%%%%%%%%%%%%%

\chapter{Metaheurísticas}
\label{cha:metah}

\section{Introducción}
\label{sec:intro}

Las metaheurísticas forman una de las familias de técnicas de optimización más utilizadas para resolver problemas complejos de 
optimización combinatoria. Estos métodos no definen una secuencia de pasos para abordar un problema concreto, sino que 
proporcionan una estrategia general de búsqueda que puede adaptarse a diferentes contextos, definiendo una representación de la 
solución, una función objetivo, y un mecanismo para generar y evaluar soluciones nuevas. Esta flexibilidad permite que una misma
metaheurística pueda emplearse para resolver problemas tan diferentes como el Problema del Viajante, la planificación de tareas, 
la asignación de recursos, o el diseño de redes.

Aunque existen muchas metaheurísticas, y cada una presenta un funcionamiento diferente, todas comparten una misma estructura 
básica. El proceso comienza con una o varias soluciones iniciales, a partir de las cuales se generan nuevas soluciones candidatas
mediante distintos mecanismos de búsqueda. Cada una de estas soluciones se evalúa mediante una función objetivo, que determina 
su calidad y decide cuáles se deben conservar para continuar el proceso. Este proceso se repite de forma iterativa hasta alcanzar
un cierto criterio de parada, como un número máximo de iteraciones, un límite de tiempo, o la ausencia de mejoras durante un 
determinado período. 

\begin{figure}[H]
  \centering 
  \fbox{\includegraphics[width=0.7\linewidth]{metah_clasif.png}}
  \caption{Clasificación de las metaheurísticas.}
  \label{fig:metah_clasif}
\end{figure}

Las principales diferencias entre unas metaheurísticas y otras residen en la forma de explorar el espacio de búsqueda, y en los
mecanismos utilizados para evitar caer en óptimos locales. Algunas incorporan elementos probabilísticos o estructuras de memoria, 
mientras que otras se inspiran en procesos biológicos o en el comportamiento colectivo observado en la naturaleza. Aun así, todas
tienen como objetivo obtener soluciones de alta calidad en un tiempo de ejecución razonable.

En las próximas secciones se van a desarrollar las metaheurísticas empleadas en este trabajo, describiendo tanto sus fundamentos
teóricos, como los aspectos necesarios para su posterior implementación y evaluación experimental.



%%% Local Variables:
%%% TeX-master: "../book"
%%% End:
